\documentclass[12pt]{article}
\usepackage{float}
\usepackage{sbc-template}
\usepackage{graphicx,url}
\usepackage[utf8]{inputenc}
\usepackage[brazil]{babel}
\usepackage{booktabs}
\usepackage{amssymb}
\usepackage[ruled,vlined]{algorithm2e}
\usepackage{amsmath}
\usepackage{pdfpages}
\usepackage[numbers]{natbib}
\usepackage{caption}
\usepackage{tabularx}
\usepackage{array}
\usepackage{subfig}

\newcolumntype{Y}{>{\raggedright\arraybackslash}X}
\sloppy

\title{Filo-Transformer: Um modelo baseado em Grafo de Alinhamento de Árvores Filogenéticas e Transformers para Identificação de Rumores e Fake News}

\author{Acauan C. Ribeiro\inst{1,2}, Eduardo L. Feitosa\inst{1}, André Carvalho\inst{1}}

\address{IComp -- Universidade Federal do Amazonas (UFAM)\\ Manaus -- AM -- Brasil
\nextinstitute
DCC -- Universidade Federal de Roraima (UFRR)\\ Boa Vista, RR -- Brasil
\email{acauan.ribeiro@ufrr.br, efeitosa@icomp.ufam.edu.br}
\email{andre@icomp.ufam.edu.br}
}

\begin{document}

\maketitle

\begin{abstract}
This paper presents Filo-Transformer, an innovative approach that combines the deep semantics of \textit{Transformer} models with the evolutionary analysis of \textit{Tree Alignment Graphs} (TAGs) to detect \textit{rumors} and \textit{fake news}. The model uses \textit{embeddings} (SBERT/GPT) to represent content and builds TAGs to trace the evolution of information, extracting attributes such as cascade depth and recombination. A Feature Tokenizer Transformer (FT-Transformer) integrates this information for classification. Experiments on the \textit{PHEME} dataset show that Filo-Transformer outperforms purely semantic models across all major metrics, with ablation analysis confirming the value of phylogenetic attributes, positioning it as an original and validated contribution to combating misinformation.
\end{abstract}

\begin{resumo}
Este artigo apresenta Filo-Transformer, uma abordagem inovadora que une a semântica profunda de modelos \textit{Transformer} com a análise evolutiva de \textit{Grafos de Alinhamento de Árvores Filogenéticas} (TAGs) para detectar \textit{rumores} e \textit{fake news}. O modelo utiliza \textit{embeddings} (SBERT/GPT) para representar o conteúdo e construir \textit{TAGs} para rastrear a evolução da informação, extraindo atributos como profundidade de cascata e recombinação. Um Feature Tokenizer Transformer (FT-Transformer) integra essas informações para classificação. Experimentos no \textit{dataset} PHEME mostram que o Filo-Transformer supera modelos apenas semânticos em todas as métricas principais, com a análise de ablação confirmando o valor dos atributos filogenéticos, oferecendo-se como uma contribuição original e validada para mitigar a desinformação.
\end{resumo}

\section{Introdução} \label{sec:introducao}
A proliferação de \textit{notícias falsas} (\textit{fake news}) em plataformas online tornou-se uma ameaça global, com implicações sociais, políticas e de segurança. Redes sociais permitem que informações não verificadas se espalhem de forma viral devido à sua natureza aberta e não moderada \cite{Zubiaga2016}. Estudos recentes revelam que \textit{notícias falsas} propagam-se mais rápido e mais longe do que notícias verdadeiras em mídias sociais. Em particular, \cite{Vosoughi2018} observaram que, no Twitter, boatos falsos difundem-se significativamente mais fundo e mais amplamente do que fatos verídicos, frequentemente por uma ordem de grandeza de diferença. Esse alcance acelerado deve-se principalmente ao compartilhamento humano e não a bots, intensificando o desafio de conter a desinformação na fonte \cite{Vosoughi2018}.

Modelos estritamente semânticos, como Bidirectional Encoder Representations from Transformers (BERT) \cite{devlin2019bert} e Sentence-BERT Adapted (SBERTA) \cite{reimers2019sentence}, são eficazes na interpretação do significado textual em um ponto específico no tempo. No entanto, esses modelos falham em capturar a dinâmica evolutiva das mensagens, negligenciando processos de reescrita, recombinação e republicação de conteúdo nas redes sociais.

Essa ausência de perspectiva temporal mascara variações textuais sutis, frequentemente cruciais para a viralização de rumores, comprometendo a capacidade desses sistemas para a detecção precoce de narrativas adulteradas em circulação.

Outro problema diz respeito à propagação, momento em que uma publicação inicial desencadeia uma série de respostas, compartilhamentos e modificações, denominada neste artigo como \textit{cascata de informação}\footnote{Ela representa o conjunto de todas as publicações inter-relacionadas que se originam de um conteúdo semente e traçam seu percurso através da rede ou plataforma.}. Os trabalhos de \cite{peng2024infodemiology, qiu2022entropy} ressaltam a importância de entender a evolução temporal e estrutural da informação, aspecto que a abordagem proposta sistematiza e explora computacionalmente.

Nesse sentido, e com o intuito de superar as limitações das abordagens existentes, este artigo investiga como integrar, eficazmente, a análise da evolução e mutação da informação com representações semânticas profundas para aprimorar a detecção de \textit{rumores} e \textit{fake news}. Destaca-se que tal integração mostra-se fundamental, uma vez que as metodologias atuais descritas na literatura frequentemente se dedicam a aspectos isolados do problema.

O objetivo geral deste trabalho é propor e avaliar um modelo híbrido inovador para detecção de \textit{rumores} e \textit{fake news}, denominado \textbf{Filo-Transformer}, que integra aspectos evolutivos e semânticos na análise textual. Especificamente, o trabalho propõe: (i) desenvolver métodos para a construção eficaz de \textit{Grafos de Alinhamento de Árvores} (TAGs) \cite{smith2013tag} a partir de \textit{cascatas informacionais} utilizando \textit{embeddings} semânticos avançados; (ii) definir e extrair atributos filogenéticos relevantes desses grafos, como profundidade e recombinação, e integrá-los com representações semânticas através de uma arquitetura Feature Tokenizer Transformer (FT-Transformer) \cite{gorishniy2021fttransformer} adaptada; (iii) e validar experimentalmente a proposta em \textit{dataset} público, comparando-a a modelos tradicionais e avaliando o impacto dos componentes filogenéticos.

Este artigo está estruturado da seguinte forma: a Seção~\ref{sec:rev-literatura} apresenta uma revisão da literatura, discutindo conceitos fundamentais e trabalhos relacionados. A Seção~\ref{sec:modelo} detalha o modelo Filo-Transformer proposto. A Seção~\ref{sec:implementacao} descreve a configuração experimental. A Seção~\ref{subsec:resultados-analise} apresenta os resultados obtidos e uma análise comparativa. A Seção~\ref{sec:trabalhos-relacionados} discute trabalhos relacionados que utilizam bases de dados e métricas similares. Finalmente, a Seção~\ref{sec:conclusao} conclui o trabalho, sintetizando as contribuições, limitações e direções para pesquisas futuras.

\section{Conceitos Fundamentais} \label{sec:rev-literatura}
Esta seção estabelece a base conceitual sobre \textit{rumores} e \textit{fake news}, e apresenta os principais conceitos e temas empregados neste artigo.

\subsection{Rumores e Fake News}
\textit{Rumores} são informações não verificadas ou incertas que circulam amplamente, frequentemente em situações onde há alta ambiguidade ou relevância social, podendo ou não ser confirmados posteriormente como verdadeiros ou falsos \cite{vosoughi2018spread}. Embora os \textit{rumores} possam surgir espontaneamente e sem intenção clara de causar danos, sua disseminação pode levar à desinformação e causar efeitos adversos em sociedades, especialmente durante crises ou eventos emergentes \cite{Zubiaga2016}. Em contraste, \textit{fake news} são intencionalmente fabricadas ou manipuladas com o objetivo explícito de enganar ou influenciar opiniões públicas, comportamentos ou decisões políticas \cite{lazer2018science,Shu2017}. A característica central das \textit{fake news} reside na intenção maliciosa e planejada de distorcer fatos ou contextos, resultando em narrativas que frequentemente exploram emoções negativas como medo, raiva ou indignação para maximizar sua disseminação e impacto \cite{Wardle2017}. Compreender a distinção e intersecção entre \textit{rumores} e \textit{fake news} é essencial para o desenvolvimento de métodos eficazes de identificação e mitigação dessas informações no ambiente digital.

\subsection{Geração de \textit{Embeddings}}
\textit{Embeddings} textuais são representações vetoriais densas e contínuas que codificam significados semânticos e sintáticos das palavras, frases ou documentos em espaços vetoriais de alta dimensão \cite{Mikolov2013}. Inicialmente, técnicas como Word2Vec e GloVe estabeleceram os fundamentos dos \textit{embeddings} baseados em redes neurais, capturando relações lineares entre palavras por meio do contexto local \cite{Pennington2014}. Recentemente, o avanço na arquitetura de modelos como o Sentence-BERT (SBERT) e os \textit{embeddings} gerados por grandes modelos de linguagem (LLMs), como GPT-3.5 e GPT-4, ampliaram significativamente a capacidade de captura semântica desses \textit{embeddings} \cite{Reimers2019,Brown2020}. Esses modelos, treinados sobre corpos de dados grandes e diversificados, conseguem produzir \textit{embeddings} altamente sensíveis às nuances semânticas, permitindo aplicações avançadas em detecção de \textit{fake news} e análise de texto \cite{devlin2019bert}.

\subsection{Transformers}
A arquitetura \textit{Transformer} representa uma mudança de paradigma em aprendizado profundo para processamento de linguagem natural (PLN), substituindo modelos recorrentes tradicionais por mecanismos baseados exclusivamente em atenção \cite{Vaswani2017}. Este mecanismo permite capturar dependências de longo alcance dentro do texto, reduzindo os problemas de desaparecimento e explosão de gradiente, característicos das redes neurais recorrentes (RNNs). \textit{Transformers} utilizam camadas de \textit{multi-head attention} que simultaneamente realizam múltiplas projeções lineares dos dados de entrada, proporcionando uma visão mais rica e diversificada das relações entre palavras ou subunidades textuais. Essa arquitetura fundamenta modelos atuais como o GPT, BERT e variantes subsequentes, os quais demonstram desempenhos excepcionais em tarefas complexas, incluindo classificação textual, análise semântica e detecção automática de \textit{fake news} \cite{Radford2019}.

\subsection{Reconstrução Filogenética Textual}
\textit{Grafos de Alinhamento de Árvores} (TAGs, do inglês \textit{Tree Alignment Graphs}) são estruturas que representam e alinham múltiplas árvores filogenéticas, capturando relações evolutivas complexas, como recombinações e divergências, originalmente em contextos biológicos \cite{smith2013tag}. Recentemente, essas técnicas têm sido aplicadas com sucesso em contextos digitais para estudar como narrativas falsas se propagam e evoluem ao longo do tempo, modelando a ancestralidade e mutação de textos por meio de similaridades calculadas com \textit{embeddings} semânticos. Tais estruturas são fundamentais para entender mecanismos subjacentes à propagação das \textit{fake news}, fornecendo uma base analítica robusta para intervenções e estratégias de mitigação no ambiente digital \cite{Jang2018}.

\section{Modelo Proposto: Filo-Transformer} \label{sec:modelo}
O Filo-Transformer é um modelo híbrido projetado para identificar \textit{rumores} e \textit{fake news}, integrando representações semânticas profundas com características extraídas da análise da evolução e propagação da informação.

\subsection{Arquitetura}
A arquitetura geral do Filo-Transformer é ilustrada na Figura \ref{fig:Modelo-Filo_T}.

\begin{figure}[!htbp]
\centering
\includegraphics[width=1\textwidth]{figures/arquitetura_filo_tranformer.pdf}
\caption{Fluxo completo do \textbf{Filo-Transformer}.}
\label{fig:Modelo-Filo_T}
\end{figure}

O método proposto compreende seis estágios principais: (1) \textbf{Entrada}, onde é realizado um pré-processamento nos dados textuais (limpeza e tokenização); (2) \textbf{Geração de \textit{Embeddings} Semânticos}, onde cada documento é convertido em um \textit{embedding} denso, denotado como $e_{\text{source}}$ para o texto de origem, ou $e_i$ e $e_j$ para representar os \textit{embeddings} de postagens específicas $p_i$ e $p_j$, normalizados por $L_{2}$; (3) \textbf{Construção de \textit{Grafos de Alinhamento de Árvores} (TAGs) e Extração de Atributos Filogenéticos}, em que os \textit{embeddings} alimentam o módulo de \emph{Reconstrução Filogenética}: calcula-se a matriz de similaridade, constroem-se grafos $k$-Nearest Neighbors (k-NN) ponderados e gera-se o \textit{Tree Alignment Graph} (TAG); (4) \textbf{Integração e Projeção dos Vetores}, onde, a partir do TAG, extraem-se atributos evolutivos, representados pelo vetor $f$, que inclui métricas como profundidade, centralidades, recombinação, entropia, \textit{embeddings} de grafo e outros que, após normalização, são projetados e concatenados aos \textit{embeddings} semânticos em um esquema \emph{dual-input}; (5) \textbf{Aprendizado Profundo com FT-Transformer}, que processa os dois vetores por camadas de \textit{multi-head attention} e redes \textit{feed-forward}, produzindo uma representação global via \textit{pooling}; e (6) \textbf{Classificação com um Neurônio de Ativação Sigmoide}, que converte essa representação na probabilidade binária: \textit{fake} (1) ou \textit{real} (0). A seguir, detalhamos melhor estas etapas.

\subsubsection{Pré-processamento e Geração de \textit{Embeddings} Semânticos}
O processamento inicial do Filo-Transformer tem início com a definição do conjunto de textos de entrada, composto por postagens potencialmente relacionadas a \textit{cascatas de desinformação}. Cada texto é submetido a um procedimento de pré-processamento automatizado, implementado com bibliotecas padrão como \texttt{nltk} e \texttt{re} em Python, que inclui: (i) remoção de URLs, menções a usuários (e.g., ``@usuário''), hashtags (e.g., ``\#tópico'') e caracteres especiais (e.g., emojis, símbolos não alfanuméricos); (ii) conversão de todos os caracteres para minúsculas; (iii) eliminação de espaços em branco redundantes; e (iv) uniformização da pontuação, substituindo variações por formas padrão (e.g., múltiplos pontos de exclamação por um único). Essas etapas são essenciais para garantir que o conteúdo analisado esteja livre de artefatos que poderiam comprometer a qualidade das representações semânticas extraídas.

Concluído o pré-processamento, cada postagem é convertida em uma representação vetorial densa, capaz de capturar as nuances semânticas e contextuais do texto. Diversos modelos de \textit{embeddings} reconhecidos na literatura foram avaliados, incluindo o \texttt{\small all-mpnet-base-v2} \cite{reimers2019sbert}, a fim de identificar a arquitetura mais adequada para a tarefa de detecção de \textit{desinformação}. Os experimentos indicaram que o modelo \texttt{\small text-embedding-3-large} da OpenAI proporcionou um ganho significativo de desempenho, com um aumento de cerca de 0.4 (ou 4 pontos percentuais) na acurácia em validações preliminares quando comparado aos demais métodos testados. Esse avanço é atribuído à maior capacidade de modelagem e ao treinamento extensivo do modelo, atualmente com 3072 dimensões, o que permite a geração de \textit{embeddings} mais ricos e informados pelo contexto. Tais representações, ao comporem a entrada semântica do Filo-Transformer, tornam-se fundamentais para diferenciar com precisão narrativas verdadeiras de conteúdos desinformativos, oferecendo uma base robusta para as etapas evolutivas e classificatórias subsequentes do modelo.

\subsubsection{Construção de Grafo de Alinhamento de Árvores (TAGs)} \label{subsec:construcao-tags}
De posse dos \textit{embeddings} semânticos dos posts a serem analisados, o próximo estágio é modelar a estrutura evolutiva da informação. Em vez de assumir uma única árvore de propagação simples, propõe-se a construção de \textit{Grafos de Alinhamento de Árvores} (TAGs). Conforme descrito na Seção~\ref{sec:rev-literatura}, os TAGs permitem modelar cenários complexos onde diferentes vertentes de uma narrativa evoluem em paralelo ou onde há incerteza sobre as relações de ancestralidade, capturando recombinações e divergências de conteúdo \cite{smith2013tag}.

A construção do TAG para uma \textit{cascata de informação} envolve os seguintes passos:
\begin{enumerate}
    \item \textbf{Identificação de Posts Semente:} Identificar os \textit{posts} iniciais que deram origem às principais narrativas dentro da cascata. Isso pode ser feito com base em temporalidade ou por agrupamento dos \textit{embeddings} semânticos.
    \item \textbf{Construção de Árvores de Propagação/Mutação:} Para cada \textit{post} semente, ou para a cascata como um todo, constroem-se uma ou mais árvores filogenéticas. Os nós da árvore são os \textit{posts}. As arestas podem ser inferidas a partir de:
    \begin{itemize}
        \item \textbf{Estrutura de Resposta:} Se a plataforma fornecer relações de conversa/respostas (e.g., \textit{replies} no Twitter/X), essa é uma forte evidência de descendência direta.
        \item \textbf{Similaridade Semântica e Temporalidade:} Um \textit{post} $p_j$ é considerado um provável descendente de $p_i$ se $p_j$ ocorreu após $p_i$ e seu \textit{embedding} $e_j$ apresenta alta similaridade semântica com $e_i$, calculada por meio de um limiar de similaridade de cosseno. O ``pai'' mais provável é o \textit{post} anterior com maior similaridade.
        \item A ``mutação'' é representada pela distância semântica entre $e_i$ e $e_j$.
    \end{itemize}
    \item \textbf{Alinhamento em um TAG:} As árvores individuais (ou sub-árvores representando diferentes linhagens da narrativa) são então alinhadas em um único TAG. Este grafo permite representar pontos de consenso, conflito (diferentes versões evoluindo de um ancestral comum) e recombinação (uma nova versão do \textit{post} que parece herdar semanticamente de múltiplos ``pais'' de diferentes ramos). A estrutura do TAG é crucial para capturar a complexidade da evolução da \textit{desinformação}, que raramente é linear. A implementação pode adaptar algoritmos de construção de TAGs de domínios biológicos, usando a similaridade semântica como análogo à similaridade genética e a estrutura temporal para guiar as relações de ancestralidade.
\end{enumerate}

\subsubsection{Extração de Atributos Filogenéticos} \label{subsec:extracao-atributos}
A seleção dos atributos filogenéticos foi guiada por hipóteses consolidadas na literatura, que indicam que \textit{rumores} e notícias verificadas exibem dinâmicas distintas em termos de evolução, profundidade e padrões de recombinação \cite{vosoughi2018spread, peng2024infodemiology}. Esses atributos, descritos na Tabela~\ref{tab:atributos_grupos}, foram definidos para quantificar diferentes facetas da propagação informacional, incluindo centralidade, estrutura local, relações ancestrais, diversidade e evolução, permitindo testar a hipótese central de que a integração de características evolutivas aprimora a detecção baseada exclusivamente em representações semânticas. Além disso, a função de perda de entropia cruzada binária foi escolhida por ser o padrão consolidado para tarefas de classificação binária, como a detecção de \textit{rumores} e \textit{fake news}, devido à sua capacidade de otimizar a separação entre classes \textit{fake} (1) e \textit{real} (0) \cite{Goodfellow-et-al-2016}. Uma vez construído o TAG, extraem-se um conjunto diversificado de atributos quantitativos, definidos para caracterizar a estrutura evolutiva e a dinâmica de propagação da informação ao longo da cascata.

\begin{table}[htbp]
\centering
\caption{Resumo dos Grupos de Atributos Extraídos do TAG}
\label{tab:atributos_grupos}
\renewcommand{\arraystretch}{1.08}
\small
\begin{tabularx}{\linewidth}{>{\bfseries}l X}
\toprule
Grupo & Atributos \\
\midrule
Centralidade \& Importância & \textit{pagerank}, \textit{closeness}, \textit{betweenness} \\
Grau \& Estrutura Local & \textit{degree\_normal}, \textit{degree\_in}, \textit{degree\_out}, \textit{is\_leaf} \\
Relações Ancestrais/Descendentes & \textit{n\_ancentrais}, \textit{n\_descendentes}, \textit{subtree\_size} \\
Diversidade \& Comunidade & \textit{gini\_similarity}, \textit{num\_comunidades} \\
Evolução, Profundidade e Recombinação & \textit{depth\_normal}, \textit{recomb\_degree}, \textit{entropy\_ancentrais}, \textit{mutacao\_rate} \\
\bottomrule
\end{tabularx}
\end{table}

Os atributos relacionados à centralidade e importância dos nós, como \textit{pagerank}, \textit{closeness} e \textit{betweenness}, têm papel fundamental para identificar pontos-chave de disseminação e possíveis alvos para controle da propagação. Esses indicadores avaliam o quão estratégico é cada nó dentro da estrutura do TAG, permitindo distinguir, por exemplo, os responsáveis pelo início de grandes cadeias de compartilhamento de informação.

No âmbito do grau e da estrutura local, valores como \textit{degree\_normal}, \textit{degree\_in}, \textit{degree\_out} e \textit{is\_leaf} oferecem uma caracterização detalhada da conectividade dos nós. Esses atributos descrevem desde o papel do nó como iniciador, intermediário ou terminal na cascata, até padrões de ramificação que podem indicar tentativas coordenadas de amplificação ou, ao contrário, trajetórias isoladas e autônomas.

A linhagem da propagação é quantificada por atributos que expressam relações ancestrais e descendentes, como \textit{n\_ancentrais}, \textit{n\_descendentes} e \textit{subtree\_size}. Esses indicadores permitem rastrear a profundidade e o alcance das narrativas, oferecendo subsídios para entender o potencial de viralização e a formação de subestruturas informacionais dentro da cascata.

Adicionalmente, atributos ligados à diversidade e à comunidade, como \textit{gini\_similarity} e \textit{num\_comunidades}, são essenciais para capturar a heterogeneidade da propagação. Eles revelam tanto a mistura de diferentes fontes de informação quanto a presença de zonas de recombinação, onde múltiplas linhagens convergem para gerar variantes inéditas do conteúdo analisado.

Por fim, as dimensões relacionadas à evolução e à recombinação da informação são modeladas por indicadores como \textit{depth\_normal}, \textit{recomb\_degree}, \textit{entropy\_ancentrais} e \textit{mutacao\_rate}. Essas medidas quantificam o grau de inovação dos nós em relação aos seus ancestrais, a profundidade evolutiva dentro da cascata e a frequência com que informações são recombinadas ou modificadas ao longo do tempo. Dessa forma, tais atributos capturam dinâmicas cruciais para a diferenciação entre o comportamento típico de notícias legítimas e os padrões característicos da disseminação de \textit{fake news} e \textit{rumores}.

\subsubsection{Integração com FT-Transformer} \label{subsec:integracao-fttransformer}
A construção de um TAG pressupõe a existência de um conjunto de publicações semanticamente correlacionadas, como variações sobre um mesmo \textit{rumor} ou narrativa específica. Para garantir coerência evolutiva e possibilitar a extração de atributos informativos, esses conjuntos devem ser previamente agrupados por tópico ou evento e, idealmente, rotulados quanto à veracidade (\textit{fake} ou \textit{real}). Essa estrutura temática comum é fundamental para que os nós da árvore compartilhem contexto suficiente, permitindo inferências válidas sobre ancestralidade, mutações semânticas e padrões de disseminação.

O Feature Tokenizer Transformer (FT-Transformer) é uma arquitetura baseada em \textit{Transformer} projetada especificamente para dados tabulares, onde cada característica (coluna) é tratada como um \textit{token}. Ele aplica o mecanismo de \textit{auto-atenção} entre as características, permitindo ao modelo aprender interações complexas entre elas. No caso proposto, as ``características'' são os componentes do vetor de \textit{embedding} semântico ($e_{\text{source}}$) e os diferentes atributos filogenéticos extraídos (vetor $f$).

\begin{enumerate}
    \item \textbf{Tokenização de Características:} Tanto as características semânticas (elementos do vetor $e_{\text{source}}$) quanto as características filogenéticas (elementos do vetor $f$) são linearmente projetadas em \textit{embeddings} de características.
    \item \textbf{Transformer Encoder:} Uma pilha de blocos \textit{Transformer} processa esses \textit{embeddings} de características. A \textit{auto-atenção} permite que o modelo pese a importância relativa de diferentes atributos semânticos e filogenéticos, e suas interações, para a tarefa de classificação.
    \item \textbf{Classificação:} A saída do \textit{Transformer}, após o processamento por \textit{pooling} global, é passada por uma camada de projeção linear que mapeia a representação combinada (de dimensão interna 192) para um único valor escalar. Este valor é então processado por uma função de ativação sigmoide, que produz uma probabilidade binária no intervalo [0,1], onde valores próximos a 1 indicam a classe \textit{fake} (\textit{rumor}/\textit{fake news}) e valores próximos a 0 indicam a classe \textit{real} (não \textit{rumor}). Para cenários de múltiplas classes, uma função \textit{softmax} pode ser utilizada, embora neste trabalho a tarefa seja estritamente binária.
\end{enumerate}

A função de perda utilizada para o treinamento do FT-Transformer é a entropia cruzada binária, adequada para tarefas de classificação binária como detecção de \textit{rumores} ou \textit{fake news}.

\subsection{Originalidade e Inovação} \label{subsec:originalidade-inovacao}
A originalidade do Filo-Transformer reside na combinação inédita de métodos evolutivos e semânticos, projetada a priori para testar a hipótese de que a integração de atributos filogenéticos com representações semânticas aprimora a detecção de \textit{desinformação}. Essa abordagem permite modelar não apenas ``o que'' é dito, mas ``como'' a narrativa evolui e se ramifica, capturando padrões evolutivos característicos de \textit{rumores} e \textit{fake news}.

Primeiramente, a pesquisa representa uma das primeiras iniciativas formais a adaptar \textit{Grafos de Alinhamento de Árvores} (TAGs), tradicionalmente utilizados em bioinformática, para a análise da evolução textual em contextos de \textit{fake news}. Essa abordagem permite modelar interações evolutivas complexas entre múltiplas versões de narrativas, capturando recombinações e divergências de conteúdo de maneira estruturada.

Além disso, o modelo propõe a extração de atributos híbridos diretamente desses grafos, integrando características estruturais da disseminação informacional com medidas da variação semântica entre os textos. Essa fusão entre estrutura e semântica permite uma análise mais profunda da dinâmica evolutiva dos \textit{rumores} e \textit{notícias falsas}.

Por fim, destaca-se a integração avançada desses atributos com \textit{embeddings} semânticos gerados por modelos de linguagem de última geração, por meio de uma arquitetura baseada no FT-Transformer. Essa integração capacita o modelo a identificar, de forma adaptativa, as dimensões mais relevantes para a predição da veracidade das narrativas, maximizando o aproveitamento das informações contidas tanto na estrutura evolutiva quanto no conteúdo semântico dos textos analisados.

\section{Implementação e Ambiente Experimental} \label{sec:implementacao}
Esta seção detalha o \textit{pipeline} experimental, as métricas de avaliação, os \textit{datasets} utilizados e os resultados comparativos do Filo-Transformer. A avaliação foca em demonstrar a eficácia do modelo proposto e o impacto positivo da incorporação de atributos filogenéticos.

\subsection{Conjunto de Dados \textit{PHEME}} \label{subsec:pheme}
Neste trabalho, utiliza-se o \textit{PHEME} \cite{Zubiaga2016}, um conjunto de dados amplamente utilizado em pesquisas sobre detecção de \textit{rumores} e \textit{fake news} em redes sociais, especialmente no contexto de eventos noticiosos e situações de crise. O \textit{PHEME} é um corpus de conversas do Twitter anotadas manualmente como \textit{rumour} ou \textit{non-rumour} em cinco eventos de \textit{breaking news} (\textit{Charlie Hebdo}, \textit{Ferguson Unrest}, \textit{Germanwings Crash}, \textit{Ottawa Shooting} e \textit{Sydney Siege}).

A Tabela~\ref{tab:pheme_stats} resume a distribuição de classes.

\begin{table}[!htp]
\small
\setlength{\tabcolsep}{6pt}
\caption{Distribuição de \textit{rumores} e \textit{não-rumores} no \textit{PHEME}.}
\label{tab:pheme_stats}
\begin{tabularx}{\linewidth}{@{}l*{3}{>{\raggedleft\arraybackslash}X}@{}}
\toprule
\textbf{Evento} & \textbf{Rumor} & \textbf{Não Rumor} & \textbf{Total} \\
\midrule
Charlie Hebdo      & 458 & 1\,621 & 2\,079 \\
Ferguson Unrest    & 284 &   859 & 1\,143 \\
Germanwings Crash  & 238 &   231 &   469 \\
Ottawa Shooting    & 470 &   420 &   890 \\
Sydney Siege       & 522 &   699 & 1\,221 \\
\addlinespace[2pt]\midrule
\textbf{Total}     & \textbf{1\,972} & \textbf{3\,830} & \textbf{5\,802} \\
\bottomrule
\end{tabularx}
\end{table}

No total, o corpus contém \textbf{5\,802} \textit{threads}, sendo aproximadamente 34\,\% \textit{rumores} e 66\,\% \textit{não-rumores}. O conjunto de dados está disponível em seu repositório oficial\footnote{\url{https://figshare.com/articles/dataset/PHEME_dataset_of_rumours_and_non-rumours/4010619?file=6453753}}.

Vale destacar que, para aumentar a \emph{generalização} do método proposto, todas as \textit{threads} dos cinco eventos foram agregadas em um único conjunto de dados, sem distinção de tópico. No pré-processamento, o rótulo textual \textit{rumour} foi convertido para \textbf{1} e \textit{non-rumour} foi convertido para \textbf{0}. Dessa forma, os modelos foram treinados e avaliados sobre $n = 5\,802$ documentos binariamente rotulados, permitindo medir o desempenho em um cenário multi-evento de \textit{rumores} e notícias genuínas.

\subsection{Pipeline Experimental e Configuração} \label{subsec:pipeline-configuracao}
O \textit{pipeline} tem início com o pré-processamento e a tokenização dos textos. Em seguida, cada \textit{tweet} é codificado pelo modelo \texttt{\textit{text-embedding-3-large}}, gerando representações densas sob a forma de \textit{embeddings} semânticos. Na terceira etapa, realiza-se a reconstrução filogenética: para cada \textit{cascata}, é construído um \textit{Tree Alignment Graph} (TAG) a partir da matriz de similaridade entre os \textit{embeddings}. Os atributos evolutivos extraídos dessas estruturas são detalhados na Seção~\ref{subsec:extracao-atributos}.

Para a avaliação dos modelos, adotou-se o procedimento de validação cruzada com cinco \textit{folds} estratificados, garantindo a preservação da proporção de classes em cada subdivisão do conjunto de dados. Dois modelos principais foram avaliados: o Filo-Transformer, que integra \textit{embeddings} semânticos e atributos filogenéticos em uma arquitetura FT-Transformer de dupla entrada, e o \textit{Baseline} Semântico, que utiliza apenas os \textit{embeddings} semânticos, servindo como referência para quantificar o ganho proporcionado pela inclusão dos atributos filogenéticos. Os resultados obtidos também foram comparados com trabalhos anteriores que utilizaram a mesma base de dados, conforme detalhado na Seção~\ref{sec:trabalhos-relacionados}.

A configuração dos hiperparâmetros contemplou o uso de três blocos \textit{Transformer}, cada um equipado com mecanismos de \textit{atenção} de quatro cabeças e dimensão interna de 192. O treinamento foi conduzido com o otimizador Adam, utilizando uma taxa de aprendizado inicial de $10^{-4}$. Foram empregados mecanismos de \textit{early stopping}, com paciência igual a 10 e monitoramento do valor de $val\ auc$, e \textit{reduce LR on plateau}, com paciência igual a 3 e fator de redução de 0,2. Cada experimento foi executado por até 100 épocas, utilizando \textit{batch size} de 64.

\subsection{Métricas de Avaliação} \label{subsec:metricas}
Os modelos foram avaliados por Acurácia, Precisão, \textit{Recall} e \textit{F1-score}, reportadas para a classe minoritária (\textit{fake news}), principal foco deste estudo.

\subsection{Implementação}
O repositório do Filo-Transformer, disponível em \url{https://github.com/filotransformer/sbseg}\footnote{\url{https://github.com/filotransformer/sbseg}}, reúne a implementação do \textit{pipeline} completo, incluindo a geração de \textit{embeddings}, a construção de grafos com \texttt{networkx} e \texttt{node2vec}, e a arquitetura do FT-Transformer. O repositório fornece ainda instruções para execução no Google Colab, facilitando a reprodução dos experimentos.

\section{Resultados e Análise} \label{subsec:resultados-analise}
A Tabela~\ref{tab:comparacao_filo_transformer} apresenta médias e desvios-padrão obtidos na validação cruzada. Em todas as métricas, o Filo-Transformer superou o \textit{baseline} semântico, corroborando a utilidade dos atributos filogenéticos na identificação de \textit{rumores}.

\begin{table}[htbp]
\centering
\small
\caption{Comparação de Métricas entre Filo-Transformer e Baseline (GPT + FT).}
\label{tab:comparacao_filo_transformer}
\begin{tabularx}{\linewidth}{>{\raggedright\arraybackslash}p{6.2cm} *{4}{>{\centering\arraybackslash}X}}
\toprule
\textbf{Modelo} & \textbf{Acurácia} & \textbf{AUC} & \textbf{Recall} & \textbf{F1-score} \\
\midrule
Baseline (GPT + FT) & 0.8871 & 0.9495 & 0.8433 & 0.8357 \\
\textbf{Filo-Transformer (GPT + Filo + FT)} & \textbf{0.8888} & 0.9489 & \textbf{0.8530} & \textbf{0.8393} \\
\bottomrule
\end{tabularx}
\end{table}

A avaliação foi realizada com \textit{cross-validation} estratificada em 5 \textit{folds}, garantindo a mesma distribuição de classes em cada partição e reduzindo a variância das estimativas. Os valores médios (Tabela \ref{tab:comparacao_filo_transformer}) mostram que o Filo-Transformer supera o \textit{baseline} em três das quatro métricas principais: acurácia (0,8888 vs 0,8871), \textbf{\textit{recall}} (0,8530 vs 0,8433) e \textit{F1-score} (0,8393 vs 0,8357). Apenas na AUC o modelo base apresenta uma vantagem marginal (0,9495 vs 0,9489). Essa diferença de 0,0006 encontra-se dentro do intervalo de confiança calculado a partir dos resultados da validação cruzada com 5 \textit{folds}, indicando que não há evidência estatística de superioridade do modelo base sobre o Filo-Transformer nesta métrica.

Como mostrado na Figura \ref{fig:exemplo_subfigure}, o ganho consistente em \textbf{\textit{recall}}, métrica que mede a capacidade de recuperar instâncias positivas (\textit{rumores}/\textit{fake news}), é especialmente relevante, pois maximizar o \textit{recall} reduz falsos-negativos e impede que \textit{notícias falsas} passem despercebidas. Em todos os cinco \textit{folds}, essa foi a métrica em que o Filo-Transformer mais se destacou, evidenciando que os atributos filogenéticos extraídos dos TAGs fornecem sinais complementares aos \textit{embeddings} semânticos. Assim, a integração de informações evolutivas torna o modelo mais robusto para detecção de \textit{desinformação}, especialmente em contextos onde omissões podem ser mais graves que alarmes falsos.

\begin{figure}[!htbp]
\centering
\includegraphics[width=0.85\textwidth]{figures/radar_new.pdf}
\caption{Comparação do modelo com e sem Filogenia}
\label{fig:exemplo_subfigure}
\end{figure}

\subsection{Análise Visual da Filogenia Textual através de Ego Graphs}
Para complementar a avaliação quantitativa do Filo-Transformer, foi realizada uma inspeção qualitativa da árvore filogenética textual (TAG) com \textit{grafos centrados} (\textit{Ego Graphs}) \cite{wasserman1994social}. Na visualização global (Figura \ref{fig:tag}), o TAG é renderizado em vermelho (\textit{rumores}) e azul claro (notícias verdadeiras), ambos com baixa opacidade, criando um fundo sutil. Em seguida, subgrafos locais centrados em um nó ego e seus vizinhos diretos (raio = 1) são destacados por variações de cor, tamanho e contorno, facilitando a interpretação de padrões estruturais; o pseudocódigo do algoritmo é apresentado ao final deste subtópico. Nós de \textit{rumores} aparecem em vermelho e notícias verdadeiras em azul, com tamanho proporcional ao grau de entrada (número de fontes pais). A borda fúcsia indica o nó central do \textit{Ego Graph}, e as arestas exibem setas na direção \textit{pai → filho}. Essa combinação de cor, escala e contorno evidencia padrões locais sem perder a percepção da estrutura global.

Foram identificados dois padrões ilustrativos, selecionados conforme as hipóteses da pesquisa:
\begin{enumerate}
    \item \textbf{\textit{Rumores} como terminais (\textit{Ego Graph 1})}: Muitos \textit{Ego Graphs} centrados em nós vermelhos apresentam quase nenhum filho, configurando-se como pontas de cadeia (``folhas''). Esse ``beco sem saída'' visual apoia a hipótese de que \textit{rumores} tendem a não gerar derivações subsequentes, sugerindo que a ausência de propagação futura é indicadora de falsidade.
    \item \textbf{Alta recombinação em \textit{rumores} (\textit{Ego Graph 2})}: Em outros subgrafos, nós vermelhos de maior diâmetro recebem várias arestas de entrada, evidenciando uma confluência de fontes. Esse mosaico de segmentos textuais reforça a hipótese de que \textit{fake news} frequentemente resultam de colagens de múltiplas origens, exibindo, portanto, maior heterogeneidade interna.
\end{enumerate}

\begin{figure}[!htbp]
\centering
\includegraphics[width=1\textwidth]{figures/ego_graph_geral.pdf}
\caption{Visão geral do TAG com \textit{Ego Graphs} destacados.}
\label{fig:tag}
\end{figure}

Embora não tenham sido destacados grafos centrados em notícias verdadeiras, sua inspeção revela cadeias de propagação mais lineares ou recombinações que preservam a veracidade, possivelmente reflexo de práticas editoriais consolidadas. Essa análise visual fornece um suporte intuitivo aos atributos estruturais que o modelo captura quantitativamente, demonstrando que a ``história evolutiva'' de um texto carrega sinais valiosos de veracidade e pode ser explorada por métodos de aprendizado profundo.

\section{Trabalhos Relacionados e Posicionamento da Proposta} \label{sec:trabalhos-relacionados}
Recentemente, diferentes propostas de detecção usando o \textit{dataset} \textit{PHEME} foram apresentadas. Esta seção sumariza essas propostas, incluindo a do presente trabalho.

\cite{sharma2024semtec} propuseram o modelo SEMTEC, que integra processamento semântico via arquiteturas \textit{transformer} com análise emocional do conteúdo textual. Essa combinação, ao focar nas características intrínsecas das mensagens, demonstrou alta acurácia ($\sim 92\%$ no \textit{PHEME}), sugerindo a importância das pistas emocionais e semânticas no processo de detecção. Por outro lado, a métrica \textit{F1-score} não foi detalhada, o que dificulta a análise em cenários de classes desbalanceadas.

\cite{li2024samgat} apresentaram a rede SAMGAT, uma arquitetura de Rede de Atenção em Grafos (GAT) que emprega mecanismos de \textit{atenção} dinâmica para diferenciar a relevância de \textit{posts} originais e comentários. O SAMGAT obteve desempenho notável no \textit{PHEME} (86,4\% acurácia, 86,3\% \textit{F1-score}), ressaltando a relevância da estrutura de interação e da \textit{atenção} em grafos para capturar o fluxo informacional. No entanto, abordagens que tratam a estrutura de propagação de forma estática ou temporalmente, tipicamente não modelam explicitamente a transformação ou evolução do conteúdo e seu significado ao longo dessa estrutura.

\cite{wu2025rumor} investigaram a dinâmica de propagação utilizando GRUs e um \textit{Temporal Tree Transformer} para codificar árvores de conversação em diferentes janelas de tempo. Embora o modelo tenha reportado métricas mais modestas (75,84\% acurácia, 71,98\% Macro \textit{F1-score} no \textit{PHEME}), os autores destacam a necessidade de considerar o aspecto temporal da disseminação, complementando as visões puramente estruturais ou de conteúdo.

A Tabela \ref{tab:comparacao-trabalhos-relacionados} sumariza o desempenho das abordagens relatadas usando o \textit{PHEME}.

\begin{table}[htb]
\centering
\caption{Comparação do Filo-Transformer com Trabalhos Relacionados.}
\label{tab:comparacao-trabalhos-relacionados}
\begin{tabular}{p{5cm}ccc}
\toprule
\textbf{Modelo} & \textbf{Acurácia} & \textbf{F1-score} & \textbf{Recall} \\
\midrule
Wu et al. (2022) – TreeRumorEval & 0.8400 & 0.8350 & 0.8100 \\
Li et al. (2024) – SAMGAT & 0.8640 & 0.8630 & 0.8420 \\
Sharma \& Srivastava (2024) – SEMTEC & 0.9200 & -- & -- \\
\textbf{Filo-Transformer (nosso)} & \textbf{0.8888} & \textbf{0.8393} & \textbf{0.8530} \\
\bottomrule
\end{tabular}
\end{table}

\subsection{Posicionamento da Filo-Transformer}
Diferentemente dos trabalhos que priorizam o conteúdo textual \cite{sharma2024semtec} ou a estrutura de propagação estática/dinâmica \cite{li2024samgat, wu2025rumor}, a proposta apresentada inaugura um novo paradigma ao \textbf{modelar explicitamente a evolução informacional e semântica} das narrativas em mídias sociais por meio da construção e análise de TAGs.

A representação via TAGs permite capturar como o conteúdo e o significado da informação se transformam e divergem na cadeia de propagação, identificando padrões evolutivos característicos de \textit{rumores} em contraste com notícias verdadeiras. Ao integrar os \textit{embeddings} semânticos de modelos de linguagem modernos com atributos extraídos dos TAGs, processados pela arquitetura FT-Transformer, o \textbf{Filo-Transformer} atinge uma compreensão mais abrangente da disseminação da informação.

Resultados no \textit{dataset} \textit{PHEME} demonstram a eficácia dessa abordagem híbrida, com 88,9\% de acurácia e AUC de 94,9\%, desempenho comparável ao SEMTEC (que não reporta \textit{F1-score}). O Filo-Transformer ainda apresenta \textit{F1-score} de 83,9\% e \textit{recall} de 85,3\%, valores próximos aos melhores resultados recentes, como o SAMGAT (86,3\% de \textit{F1-score}). Esses números ressaltam a capacidade do modelo em lidar com o desbalanceamento do \textit{PHEME}, mantendo sensibilidade na identificação da classe minoritária.

Assim, o Filo-Transformer representa um avanço relevante não só pela performance robusta, mas também pela introdução da perspectiva filogenética na modelagem da evolução semântica, abrindo novas possibilidades para pesquisas em detecção de \textit{desinformação} baseada em processos evolutivos.

\section{Conclusão}
\label{sec:conclusao}

Este artigo aborda o problema da detecção de \textit{desinformação} em ambientes digitais, propondo o \textbf{Filo-Transformer}, um modelo híbrido que integra a profundidade da análise semântica com a modelagem evolutiva da propagação informacional.

O objetivo da pesquisa é investigar como a combinação entre representações semânticas densas, \textit{Grafos de Alinhamento de Árvores} (TAGs) e arquiteturas \textit{Transformer} pode aprimorar a identificação de \textit{rumores} e \textit{fake news}.

Os objetivos definidos foram alcançados por meio do desenvolvimento de um \textit{pipeline} composto por: (1) geração de \textit{embeddings} semânticos de alta qualidade; (2) construção de TAGs que representam a propagação e mutação de narrativas, adaptando conceitos da biologia computacional ao domínio textual; (3) extração de atributos filogenéticos informativos (como profundidade, recombinação, ancestralidade e taxa de mutação semântica); e (4) fusão desses atributos com representações semânticas em um modelo Feature Tokenizer Transformer (FT-Transformer) para classificação robusta.

As principais contribuições do trabalho são: (i) o \textbf{Filo-Transformer}, um modelo híbrido que integra análise filogenética via TAGs e \textit{Transformers} para representação semântica; (ii) a \textbf{aplicação de TAGs ao domínio da \textit{desinformação}}, com formalismo para modelar a evolução de narrativas digitais, considerando múltiplas linhagens e recombinações; (iii) a \textbf{definição de novos atributos evolutivos}, que quantificam a dinâmica de mutação e propagação da informação; e (iv) uma \textbf{avaliação experimental rigorosa}, demonstrando a superioridade do Filo-Transformer sobre abordagens baseadas apenas em conteúdo semântico no \textit{PHEME}, com ganhos consistentes em acurácia, precisão, \textit{recall} e \textit{F1-score}.

\subsection{Limitações}
Apesar dos resultados promissores, o estudo apresenta limitações. A construção de TAGs a partir de dados textuais ruidosos e em larga escala pode ser computacionalmente custosa, além de depender de heurísticas para inferência de ancestralidade. A interpretabilidade dos atributos evolutivos e sua contribuição para a decisão do FT-Transformer ainda pode ser mais explorada. Ademais, os experimentos foram realizados principalmente no conjunto \textit{PHEME}, restrito a \textit{tweets} em inglês, o que limita a generalização para outros contextos linguísticos e domínios de \textit{desinformação}.

\subsection{Trabalhos Futuros}
Com base nas limitações descritas, são propostas diversas direções para investigações futuras:
\begin{itemize}
    \item \textbf{Expansão para Cenários Multilíngues e Multimodais:} Estender o modelo para outros idiomas, com uso de \textit{embeddings} multilíngues, bem como incorporar análise de conteúdo multimodal (texto, imagem, vídeo), onde a recombinação pode ocorrer em diferentes dimensões da informação.
    \item \textbf{Incorporação de Informações Temporais e Geográficas:} Integrar atributos temporais refinados (como velocidade de mutação) e dados georreferenciados, quando disponíveis e eticamente viáveis, para enriquecer a modelagem filogenética.
    \item \textbf{Otimização da Construção de TAGs:} Investigar métodos mais eficientes e escaláveis para construção e alinhamento de TAGs, possivelmente com técnicas de aprendizado por reforço para otimizar heurísticas.
    \item \textbf{Análise de Causalidade e Intenção:} Explorar se padrões evolutivos extraídos dos TAGs podem contribuir para a inferência de intenção (maliciosa ou não) ou identificação de atores coordenadores de campanhas de \textit{desinformação}.
\end{itemize}

\noindent{\textbf{Agradecimentos}}

Agradece-se à UFRR e à UFAM pelo apoio à pesquisa, e aos professores Dr. Guilherme Pimentel Telles (IC/Unicamp) e Dra. Rosane Minghim (UCC/Irlanda) pelas ideias fundamentais que possibilitaram o desenvolvimento do trabalho.
O presente trabalho foi realizado com apoio da Coordenação de Aperfeiçoamento de Pessoal de Nível Superior - Brasil (CAPES-PROEX) - Código de Financiamento 001. Este trabalho foi parcialmente financiado pela Fundação de Amparo à Pesquisa do Estado do Amazonas – FAPEAM – por meio do projeto POSGRAD 2024/2025.

% -------------------------------------------------------------------
% -- FIM DO TEXTO ----------------------------------------------------------
% -------------------------------------------------------------------

\bibliographystyle{sbc}
\bibliography{sbc-template}
\end{document}